\section{Theoretischer Hintergrund}\label{sec:2}

In diesem Abschnitt wird in das zugrundeliegende Framework der Arbeit (Abschnitt \ref{sec:sfl}), sowie die abgedeckten Disziplinen von Linguistik und Rechtwissenschaften eingeführt.


\subsection{Systemic Functional Linguistics (SFL)}\label{sec:2:sfl}

\textit{Systemic Functional Linguistics} (kurz und im Folgenden SFL) hat ihren Ursprung in den 1970er Jahren. M.A.K. Halliday setzte seine Ideen den in den 70ern und 80ern weit verbreiteten strukturalistischen und/oder generativen Ansätzen (u.a. denen von Noam Chomsky) gegenüber und propagierte, dass Sprache ein semiotisches System bzw. ein Systemnetzwerk aus mehreren Subsystemen ist \citep{Halliday1978}, das stets vor dem Kontext seiner Funktionen interpretiert werden muss - nicht als Menge von Regeln und separat analysierbarer Einheiten.

Das Fundament für Hallidays Theorie geht zurück auf de Saussure (``Language is a social fact'', zitiert nach \citet[1]{Halliday1978}). Die zugrundeliegende Idee besteht nach Halliday darin, Sprache vor einem soziokulturellen Kontext zu interpretieren, der wiederum die Kultur als semiotisches System umfasst \citep[2]{Halliday1978}. Somit liegt bereits hier ein Fokus auf dem semantischen Beitrag soziokultureller Aspekte (bzw. genauer: Zeichen und deren Zusammensetzung; im Sinne der Semiotik).

%Halliday wiederum passte diesen im Sinne von J.R. Firth an. Nach Firth \todo{quelle} ist Sprache zwar ein System aus einer Menge von Formen und damit korrelierenden Bedeutungen. Jedoch funktioniert Sprache nicht über eindeutige Zuweisungen der Werte zueinander. Vielmehr besteht jedes System aus einer Menge von Optionen bezüglich der Form oder der Bedeutung eines Ausdruckes. Bedeutungs- und Formunterschiede entstehen durch unterschiedliche Entscheidungen hinsichtlich der Optionen. Allerdings umfassen die Optionen nie den gesamten Bedeutungsbeitrag oder die gesamte Form, sondern nur ein Merkmal.

\subsubsection{Struktur der SFL}

Die SFL strukturiert Sprache als System mit zwei Ebenen; Ausdruck -- \textit{expression} -- und Inhalt -- \textit{content} \citep[24]{Halliday2004}. Die Inhaltsebene enthält zwei weitere; Lexikogrammatik und Semantik. Die Ausdrucksebene enthält Phonetik und Phonologie einer Sprache und soll aufgrund der Datengrundlage schriftsprachlicher Texte hier nicht weiter behandelt werden. Insbesondere sind sprachliche Äußerungen jedweder Art in der SFL als \textit{Text} zu bezeichnen; auch, wenn diese gesprochen und nicht geschrieben vorliegen \citep[26]{Halliday2004}.

Der Systembegriff in der SFL ist auf zwei verschiedene Weisen zu verstehen. Einerseits als Menge von Subsystemen und Funktionsweisen einer Sprache, die Text produziert, andererseits als das zugrundeliegende Potential einer Sprache zur Konstruktion von Bedeutung (ebd.).
Letzterer bildet gemeinsam mit dem Textbegriff ein Kontinuum und ist gleichzeitig eng verwand mit dem Registerbegriff. Das System stellt funktionale Variation dar, genau wie Hallidays Registerdefinition \citep[27]{Halliday2004}, und diese wiederum ist in der SFL eine Menge von Wahrscheinlichkeiten von Variationen (z.B. die Wahrscheinlichkeit dafür, dass in Wettervorhersagen eher das Futur verwendet wird als in einem Märchen).

Sprache hat nach Halliday mehrere Metafunktionen; die ideationale (``language as reflection''), die interpersonelle (``language as action) und die textuelle
\citep[29--30]{Halliday2004}. Der Term \textit{Metafunktion} geht dabei auf die Intetion zurück, ihn von dem einer \textit{Funktion} abzugrenzen - nach Halliday sind Funktionen in jeder Grammatiktheorie -- sehr allgemein formuliert -- als ''purpose or way of using language`` \citep[30]{Halliday2004}.

Um die Grammatik einer Sprache zu untersuchen, bedient sich Halliday weiterer Terminologie. Zunächst beschreiben \textit{Klassen} im konservativen Sinne je eine Menge von Objekten, die gewissen Eigenschaften teilen. In der SFL ist die Klasse eines Objektes ein Indikator dafür, welche grammatikalischen Funktionen das zugehörige Objekt potentiell erfüllen kann.

Konventionelle Klassen sind u.a. der Verbkomplex als funktionales Gefüge (Verbgruppe; \textit{verbal group}) und die Nominalgruppe (\textit{nominal group}). Die Verbgruppe hat die Funktion, die Art des Vorgangs in der clause zu klassifizieren (den \textit{process type}, s.u), die Nominalgruppe klassifiziert semantische Rollen (Subject, Actor, Theme), \cite[53--58]{Halliday2004}.

\subsubsection{Funktionen der Nominalgruppe}

Diese Arbeit wird sich auf die Funktion der Verbgruppe fokussieren. Nichtsdestotrotz werden die Nomen in die Analyse miteinbezogen.

Je nachdem, welche Funktion die Nominalgruppe im Satz einnimmt, kann man es als \textit{Subject}, \textit{Actor} oder \textit{Theme} klassifizieren.

Ersteres ist der Fall, wenn man den Satz als \textit{exchange} analyisert, in dem Information zwischen Sender und Empfänger ausgetauscht wird. Halliday beschreibt das \textit{subject} als die Garantie für die Richtigkeit der ausztauschenden Information.

Die Nominalgruppe ist \textit{Theme}, wenn der Satz als \textit{message} analysiert wird. Das \textit{Theme} wird beschrieben als das \textit{grounding} dafür, was ausgesagt werden soll.

Analysiert man den Satz als \textit{representation}, nimmmt die Nominalgruppe die Funktion des \textit{Actor} ein, also als die die beschriebene Handlung ausführende Instanz.


%clause as message, as exchange, as representation

\subsubsection{Das \textit{transitivity system} in der SFL}

Wie oben beschrieben, besteht ein \textit{System} aus einer Menge von Optionen, die gewählt oder verworfen werden können, während eine sprachliche Äußerung gebildet wird. Eines dieser Systeme ist das \textit{transitivity system}, welches die funktionalen, zu wählenden Optionen für Verbgefüge beschreibt. Verben werden im Rahmen der SFL (und dieser Arbeit) nicht rein syntaktisch analysiert, sondern stets als funktionale Gefüge (\textit{verbal groups}).

Die Optionen, die für die Funktion eines Verbs zur Verfügung stehen, sind die folgenden.

\begin{table}
\centering
\begin{tabular}{ll}
\toprule
process type & Beispiel\\
\midrule
material & (z.B. blühen, anfassen)\\
mental & (z.B. fühlen, bewegen, )\\
relational & \\
behavioural & \\
verbal & \\
existential & \\
\bottomrule
\end{tabular}
\end{table}



\subsection{Korpuslinguistik}
\label{sec:2:kl}

Motivation für korpusbasierte SFL-Analysen findet sich bereits in \citet[34--35]{Halliday2004}. Durch die Anreicherung großer Sprachdatenmengen ist es nicht nur möglich, Grammatiken aufgrund einzelner Texte zu konzipieren. Vielmehr kann Grammatik auch quantitativ untersucht werden. Jedoch merkt Halliday auch an, dass die Relevanz der Theorieforschung nicht der Relevanz der Sammlung von Daten unterliegen sollte. Er plädiert gegen die anti-theoretischen Ideologien des späten 20. Jahrhunderts, weswegen in dieser Arbeit auch qualitativ -- und in diesem Zuge auch auf Basis von Einzelfällen -- gearbeitet wird. Zudem spricht er automatisierten Analysen ab, eine vollumfängliche funktionale Analyse von clauses fassen zu können, auch semantische Analysen lägen außerhalb der Möglichkeiten korpusbasierter Analysen.

Von \citet[49]{Halliday2004} getestete Tools für manuelle Analysen sind zum Beispiel \textit{SysFan}, \textit{Functional Grammar Processor} und \textit{Systemics}, für automatisierte Analysen wurde {MonoConc}, {WordSmith}, {SysConc} und \textit{COBUILD tools} getestet.

\subsection{Rechtslinguistik}
\label{sec:2:ll}

\subsection{Rechtswissenschaftlicher Kontext}

\subsubsection{Rechtsgeschichte}
\label{sec:2:rg}

\subsubsection{Verfassungsrecht}
\label{sec:2:vr}
